% GENERATED FILE. Edit canonical lessons, then run npm run generate:books.
% canonical-source-hash: fnv1a64:45eceae2ad2e2e94
% canonical-lessons: TE-C62-thread, TE-C62-needle, TE-C62-broom, TE-C62-umbrella, TE-C62-mirror

\chapter{Things People Make}
\label{ch:te-made-things}
\hlchaptermodality{62}

\begin{chapteropening}
\textbf{By the end of this chapter:} \emph{I can name five things a Telugu household makes or keeps, and ask politely for one of them.}

Complete TE-C62-mirror and demonstrate the chapter promise.
\end{chapteropening}

\section[dāraṁ]{\te{దారం} (dāraṁ) — thread}

\label{lesson:TE-C62-thread}



\begin{quote}
Before the new one: what did \emph{mirapa} mean?
\end{quote}

\subsection*{You'll want to know: \te{దారం}}
\textbf{\te{దారం}} (\emph{dāraṁ}) — \textquotedblleft{}thread\textquotedblright{}.

\te{దారం} is thread — spun cotton, fine enough to sew with and strong enough to tie a knot that holds. It is the first thing in this chapter and the smallest.

The word came in from Sanskrit, where it belongs to a family about tearing and drawing out. Thread is what you draw out of a mass of cotton.

\te{దారం} does ceremonial work too. The thread tied round a wrist at a festival is a \te{దారం}, and once it is on it stays there until it falls off by itself.

\begin{grammarlens}[title={What you've built}]
The first thing a house makes for itself.
\end{grammarlens}

\subsection*{Your turn}
Say these aloud:

\begin{itemize}
\raggedright
  \item \emph{dāraṁ}
  \item it once more, slowly
  \item \emph{dāraṁ}, then \emph{dayacēsi}, and ask for it politely
\end{itemize}

\begin{tcolorbox}[breakable,colback=teal!4,colframe=teal!35!black,title={Before you move on}]
What does \te{దారం} mean? (\textquotedblleft{}thread\textquotedblright{}.) Say it once more.
\end{tcolorbox}
\section[sūdi]{\te{సూది} (sūdi) — a needle}

\label{lesson:TE-C62-needle}



\begin{quote}
Before the new one: what did \emph{dāraṁ} mean?
\end{quote}

\subsection*{You'll want to know: \te{సూది}}
\textbf{\te{సూది}} (\emph{sūdi}) — \textquotedblleft{}a needle\textquotedblright{}.

\te{సూది} is a needle. It arrived from Sanskrit \emph{sūcī}, and Telugu wore the middle sound down to a \emph{d} — a change the language makes often enough that the two shapes line up easily once you know to expect it.

A \te{సూది} is useless without a \te{దారం}, and a \te{దారం} is only a tangle without a \te{సూది}. That is about as neat a pair as this book has come across.

The same word does the \te{వైద్యుడు}'s work: an injection is a \te{సూది}, and somebody at a clinic who uses the word is not talking about mending a \te{బట్ట}.

\begin{grammarlens}[title={What you've built}]
Two, and neither is much use without the other.
\end{grammarlens}

\subsection*{Your turn}
Say these aloud:

\begin{itemize}
\raggedright
  \item \emph{sūdi}
  \item it once more, slowly
  \item \emph{sūdi}, then \emph{dāraṁ}, and say which one goes through the other
\end{itemize}

\begin{tcolorbox}[breakable,colback=teal!4,colframe=teal!35!black,title={Before you move on}]
What does \te{సూది} mean? (\textquotedblleft{}a needle\textquotedblright{}.) Say it once more.
\end{tcolorbox}
\section[cīpuru]{\te{చీపురు} (cīpuru) — a broom}

\label{lesson:TE-C62-broom}



\begin{quote}
Before the new one: what did \emph{sūdi} mean?
\end{quote}

\subsection*{You'll want to know: \te{చీపురు}}
\textbf{\te{చీపురు}} (\emph{cīpuru}) — \textquotedblleft{}a broom\textquotedblright{}.

\te{చీపురు} is the broom: a bundle of coconut-leaf ribs bound tight at one end, used bent at the waist rather than standing upright.

The word is inherited. Sweeping is the first work of the morning in a Telugu house, done before the \te{ముగ్గు} goes down at the doorstep.

A \te{చీపురు} is treated with a little care because of that. It is not stood on its bristles, it is not stepped over, and it is not lent out lightly.

\begin{grammarlens}[title={What you've built}]
Three, and this one starts the day.
\end{grammarlens}

\subsection*{Your turn}
Say these aloud:

\begin{itemize}
\raggedright
  \item \emph{cīpuru}
  \item it once more, slowly
  \item \emph{cīpuru}, then \emph{muggu}, and say which one comes first
\end{itemize}

\begin{tcolorbox}[breakable,colback=teal!4,colframe=teal!35!black,title={Before you move on}]
What does \te{చీపురు} mean? (\textquotedblleft{}a broom\textquotedblright{}.) Say it once more.
\end{tcolorbox}
\section[goḍugu]{\te{గొడుగు} (goḍugu) — an umbrella}

\label{lesson:TE-C62-umbrella}



\begin{quote}
Before the new one: what did \emph{cīpuru} mean?
\end{quote}

\subsection*{You'll want to know: \te{గొడుగు}}
\textbf{\te{గొడుగు}} (\emph{goḍugu}) — \textquotedblleft{}an umbrella\textquotedblright{}.

\te{గొడుగు} is an umbrella, and in Telugu it answers the sun as much as the rain — which is the commoner use of it on a hot morning.

The word is inherited. Long before cloth and wire there was a \te{గొడుగు} of palm leaf, and the ceremonial one held over a deity or over a bride is still made and named the same way.

It settles two things already met: \te{గాలి} will turn it inside out and \te{మంచు} will not, so a \te{గొడుగు} goes out in one kind of weather and stays at home in the other.

\begin{grammarlens}[title={What you've built}]
Four.
\end{grammarlens}

\subsection*{Your turn}
Say these aloud:

\begin{itemize}
\raggedright
  \item \emph{goḍugu}
  \item it once more, slowly
  \item \emph{goḍugu}, then \emph{gāli}, and say what one does to the other
\end{itemize}

\begin{tcolorbox}[breakable,colback=teal!4,colframe=teal!35!black,title={Before you move on}]
What does \te{గొడుగు} mean? (\textquotedblleft{}an umbrella\textquotedblright{}.) Say it once more.
\end{tcolorbox}
\section[addaṁ]{\te{అద్దం} (addaṁ) — a mirror}

\label{lesson:TE-C62-mirror}



\begin{quote}
Before the new one: what did \emph{goḍugu} mean?
\end{quote}

\subsection*{You'll want to know: \te{అద్దం}}
\textbf{\te{అద్దం}} (\emph{addaṁ}) — \textquotedblleft{}a mirror\textquotedblright{}.

\te{అద్దం} is a mirror. Of the five things here it is the only one a household buys rather than makes, and the only one wearing a borrowed name — from Sanskrit \emph{ādarśa}, \textquotedblleft{}what shows back\textquotedblright{}.

The other four come out of whatever grows nearby: a \te{దారం} out of cotton, a \te{చీపురు} out of leaf ribs, a \te{గొడుగు} out of palm and cane, and a \te{సూది} out of a scrap of metal from the smith at the edge of the \te{గ్రామం}. An \te{అద్దం} needs a factory.

Asking for one is where \te{దయచేసి} earns its keep. \te{అద్దం} \te{ఇవ్వండి} on its own lands like an instruction; with \te{దయచేసి} in front it lands like a request.

\begin{grammarlens}[title={What you've built}]
Five: \te{దారం}, \te{సూది}, \te{చీపురు}, \te{గొడుగు}, \te{అద్దం}. Enough to ask for what you need with \te{దయచేసి} in front of it.
\end{grammarlens}

\subsection*{Your turn}
Say these aloud:

\begin{itemize}
\raggedright
  \item \emph{addaṁ}
  \item it once more, slowly
  \item all five in order, then say \emph{dayacēsi addaṁ ivvaṇḍi}
\end{itemize}

\begin{tcolorbox}[breakable,colback=teal!4,colframe=teal!35!black,title={Before you move on}]
What does \te{అద్దం} mean? (\textquotedblleft{}a mirror\textquotedblright{}.) Say it once more.
\end{tcolorbox}
